\documentclass[11pt]{article}

\usepackage[a4paper,margin=1.05in]{geometry}
\usepackage{amsmath,amssymb,amsthm,mathtools}
\usepackage{enumitem}
\usepackage[colorlinks=true,linkcolor=blue,citecolor=blue,urlcolor=blue]{hyperref}

\numberwithin{equation}{section}

\newtheorem{theorem}{Theorem}[section]
\newtheorem{lemma}[theorem]{Lemma}
\newtheorem{proposition}[theorem]{Proposition}
\newtheorem{corollary}[theorem]{Corollary}
\newtheorem{conjecture}[theorem]{Conjecture}
\theoremstyle{definition}
\newtheorem{definition}[theorem]{Definition}
\newtheorem{remark}[theorem]{Remark}

\newcommand{\R}{\mathbb{R}}
\newcommand{\St}{\operatorname{St}}
\newcommand{\tr}{\operatorname{tr}}
\newcommand{\diag}{\operatorname{diag}}
\newcommand{\spec}{\operatorname{spec}}
\newcommand{\rank}{\operatorname{rank}}
\newcommand{\conv}{\operatorname{conv}}
\newcommand{\smin}{\sigma_{\min}}
\newcommand{\lmin}{\lambda_{\min}}
\newcommand{\one}{\mathbf{1}}

\title{Notes Toward the Goreinov--Tyrtyshnikov--Zamarashkin Hypothesis in the First Open Case}
\author{Draft template for the case $(n,k)=(6,3)$}
\date{\today}

\begin{document}
\maketitle

\begin{abstract}
This is a working template for a proof attempt in the first open case of the
Goreinov--Tyrtyshnikov--Zamarashkin row-selection hypothesis.  The target is the
statement that every real $6\times 3$ matrix with orthonormal columns contains a
$3\times 3$ row submatrix whose least singular value is at least $1/\sqrt 6$.

The document deliberately separates proved reductions and structural lemmas from
conjectural or numerical material.  It records: the formulation in terms of
principal submatrices of a rank-three orthoprojector; the single-pair extension
mechanism; the equal-leverage slice and its involution model; two obstructions to
natural proof strategies; and a minimax/KKT reformulation that appears to be the
most promising remaining route.
\end{abstract}

\tableofcontents

\section{The target}

Let
\[
  \St(n,k)=\{A\in\R^{n\times k}: A^T A=I_k\}.
\]
For $A\in\St(n,k)$ write its rows as $r_1,\ldots,r_n\in\R^k$, set
\[
  P=AA^T,\qquad \ell_i=P_{ii}=\|r_i\|^2,\qquad c_{ij}=P_{ij}=\langle r_i,r_j\rangle.
\]
Thus $P$ is a rank-$k$ orthogonal projector.  For $I\subset[n]$ let $P_{II}$ be
the corresponding principal submatrix.  If $|I|=k$, then
\[
  P_{II}=A_I A_I^T,
\]
so the eigenvalues of $P_{II}$ are the squared singular values of $A_I$.

\begin{conjecture}[GTZ in the first open case]\label{conj:gtz63}
For every $A\in\St(6,3)$ there is a triple $I\subset[6]$ such that
\[
  \lmin(P_{II})\ge \frac16.
\]
Equivalently, $\smin(A_I)\ge 1/\sqrt6$ for some $3\times3$ row submatrix.
\end{conjecture}

It is useful to write
\[
  F(A)=\max_{|I|=3}\lmin(P_{II}),\qquad f(A)=\max_{|I|=3}\smin(A_I).
\]
Then $F(A)=f(A)^2$, and Conjecture~\ref{conj:gtz63} is
\[
  \min_{A\in\St(6,3)} F(A)\ge \frac16.
\]

\section{Small leverage reduction}

The standard Case A reduction works for all $k$.  In the present case it says
that rows with leverage at most $1/6$ can be removed at no loss, provided the
smaller instance is known.

\begin{proposition}[Case A reduction, specialized]\label{prop:caseA}
Let $A\in\St(6,3)$ and suppose some row has $\ell_i\le 1/6$.  Rotate columns so
that this row is $(b,0,0)$, where $b^2=\ell_i$, delete the row, and renormalize
the first column by $(1-\ell_i)^{-1/2}$.  This gives a matrix
$\widetilde B\in\St(5,3)$.  For every triple $J$ of rows of $\widetilde B$,
with corresponding triple $\varphi(J)\subset[6]\setminus\{i\}$,
\[
  \sqrt{1-\ell_i}\,\smin(\widetilde B_J)
  \le \smin(A_{\varphi(J)})
  \le \smin(\widetilde B_J).
\]
Consequently, if the GTZ bound is known for $(5,3)$, then it holds for every
$A\in\St(6,3)$ with a row of leverage at most $1/6$.
\end{proposition}

\begin{proof}
After the rotation and deletion, the remaining columns have Gram matrix
$\diag(1-\ell_i,1,1)$.  Multiplication of the first column by
$(1-\ell_i)^{-1/2}$ restores orthonormality.  Pulling a triple back multiplies
on the right by $\diag(\sqrt{1-\ell_i},1,1)$, and the displayed singular value
bounds follow from the min--max principle.  If
$\smin(\widetilde B_J)\ge 1/\sqrt5$, then
\[
  \smin(A_{\varphi(J)})
  \ge \sqrt{1-\ell_i}\,\frac1{\sqrt5}
  \ge \sqrt{\frac56}\,\frac1{\sqrt5}
  =\frac1{\sqrt6}.
\]
\end{proof}

\begin{remark}
This reduction is useful but does not by itself settle $(6,3)$, since it uses
the neighboring case $(5,3)$.  By spectral duality, $(5,3)$ is equivalent to
$(5,2)$, which is covered by the known $k=2$ theorem.  Thus the genuinely new
part of $(6,3)$ may be restricted to the core region
\[
  \ell_i>1/6\quad\text{for every }i.
\]
\end{remark}

\section{The single-pair extension mechanism}

Fix $A\in\St(6,3)$ and a pair $\{i,j\}$.  Let
\[
  \Gamma=P_{\{i,j\}\{i,j\}}
  =
  \begin{pmatrix}
    \ell_i & c_{ij}\\
    c_{ij} & \ell_j
  \end{pmatrix},
  \qquad
  \spec(\Gamma)=\{\mu_1,\mu_2\},\quad \mu_1\ge\mu_2.
\]
For $p\notin\{i,j\}$ set
\[
  v_p=(c_{ip},c_{jp})^T.
\]
Assume $\mu_2>1/6$, so $\Gamma-\frac16 I_2$ is positive definite, and define
\[
  q(p)=\ell_p-\frac16
  -v_p^T\left(\Gamma-\frac16 I_2\right)^{-1}v_p.
\]

\begin{lemma}[Schur criterion]\label{lem:schur}
For $p\notin\{i,j\}$,
\[
  \lmin(P_{\{i,j,p\},\{i,j,p\}})\ge\frac16
  \quad\Longleftrightarrow\quad
  q(p)\ge0.
\]
\end{lemma}

\begin{proof}
Subtract $\frac16 I_3$ from the triple Gram matrix and take the Schur
complement of the positive block $\Gamma-\frac16 I_2$.
\end{proof}

Define, for $\mu>1/6$,
\[
  h(\mu)=\frac{1-\mu}{6\mu-1}.
\]

\begin{lemma}[Extension lemma]\label{lem:extension}
If a pair satisfies
\[
  \mu_2>\frac16,\qquad h(\mu_1)+h(\mu_2)\le\frac13,
\]
then some triple containing that pair is $1/6$-good.
\end{lemma}

\begin{proof}
The moment identity
\[
  \sum_{p\notin\{i,j\}} v_pv_p^T=\Gamma-\Gamma^2
\]
and $\sum_p\ell_p=3$ give
\[
  \sum_{p\notin\{i,j\}} q(p)
  =
  \frac73
  -
  \frac{5\mu_1}{6\mu_1-1}
  -
  \frac{5\mu_2}{6\mu_2-1}
  =
  \frac13-h(\mu_1)-h(\mu_2).
\]
Thus the assumed inequality makes the sum of the four Schur slacks nonnegative.
At least one $q(p)$ is nonnegative, and Lemma~\ref{lem:schur} applies.
\end{proof}

\begin{corollary}[A trace branch]\label{cor:trace}
If some pair satisfies
\[
  \ell_i+\ell_j\ge\frac{13}{9},
\]
then Conjecture~\ref{conj:gtz63} holds for $A$.
\end{corollary}

\begin{proof}
Let $T=\ell_i+\ell_j=\mu_1+\mu_2$.  Since $\mu_1\le1$, the inequality
$T\ge13/9$ implies $\mu_2\ge4/9>1/6$.  The function $h$ is decreasing and
convex on $(1/6,\infty)$.  The maximum of
$h(\mu_1)+h(13/9-\mu_1)$ on the admissible interval occurs at an endpoint,
and the endpoint values give at most $1/3$.  Lemma~\ref{lem:extension}
finishes the proof.
\end{proof}

\section{The equal-leverage slice}

The slice
\[
  \ell_1=\cdots=\ell_6=\frac12
\]
has a particularly rigid form.  Put
\[
  J=2P-I_6.
\]
Then $J$ is a symmetric orthogonal involution with zero diagonal and
\[
  \spec(J)=\{+1,+1,+1,-1,-1,-1\}.
\]
Conversely, any such $J$ gives a rank-three projector $P=(I+J)/2$ with
diagonal entries $1/2$.

Write
\[
  a_{ij}=J_{ij}\quad (i\ne j).
\]
Equivalently, after setting $u_i=\sqrt2\,r_i$, the $u_i$ are unit vectors in
$\R^3$ with
\[
  \sum_{i=1}^6 u_i u_i^T=2I_3,\qquad a_{ij}=\langle u_i,u_j\rangle.
\]
The involution identity gives the useful row and flow equations
\[
  \sum_{k\ne i}a_{ik}^2=1,\qquad
  \sum_{m\notin\{i,j\}}a_{im}a_{jm}=0.
\]

\begin{lemma}[Slice criterion for a good triple]\label{lem:slicecriterion}
For a triple $T=\{i,j,k\}$ set
\[
  p_T=a_{ij}^2+a_{ik}^2+a_{jk}^2,\qquad q_T=a_{ij}a_{ik}a_{jk}.
\]
Then $T$ is $1/6$-good if and only if
\[
  q_T\ge \frac{p_T}{3}-\frac4{27}.
\]
\end{lemma}

\begin{proof}
On the slice,
\[
  P_{TT}=\frac12(I_3+J_{TT}).
\]
Thus $T$ is $1/6$-good iff $\lmin(J_{TT})\ge-2/3$.  The zero-diagonal
$3\times3$ matrix $J_{TT}$ has characteristic polynomial
\[
  x^3-p_Tx-2q_T.
\]
Since $J_{TT}$ is a principal submatrix of an involution, its eigenvalues lie
in $[-1,1]$ and have sum zero, so at most one can be below $-2/3$.  Therefore
the sign of $\det(J_{TT}+\frac23 I_3)$ detects whether the least eigenvalue is
at least $-2/3$.  But
\[
  \det(J_{TT}+\tfrac23 I_3)=\frac8{27}-\frac23p_T+2q_T,
\]
which gives the criterion.
\end{proof}

The Veronese reformulation is also useful.  Let
\[
  w_i=u_i u_i^T-\frac13I_3\in\operatorname{Sym}_0(\R^3).
\]
Then $\sum_iw_i=0$ and, for a triple $T=\{i,j,k\}$,
\[
  \tau_T=\tr(w_iw_jw_k)=q_T-\frac{p_T}{3}+\frac29.
\]
Thus Lemma~\ref{lem:slicecriterion} is equivalent to
\[
  T\text{ is good}\quad\Longleftrightarrow\quad \tau_T\ge\frac2{27}.
\]
Moreover,
\[
  \sum_{|T|=3}\tau_T=\frac49,\qquad
  \sum_{|T|=3}p_T=12,\qquad
  \sum_{|T|=3}q_T=0.
\]

\begin{remark}[Why averaging alone fails]
The mean value of $\tau_T$ over the $20$ triples is $1/45$, which is below
the required threshold $2/27$.  Therefore the slice theorem, if true, must use
more than unsigned averaging of the $\tau_T$.
\end{remark}

\section{Pinned block spectra}

The involution model also gives exact spectra for large principal blocks.
These facts are potential constraints in an active-set or algebraic proof.

\begin{lemma}[Five-by-five blocks]
Let $J$ be any symmetric involution with
\[
  \spec(J)=\{+1,+1,+1,-1,-1,-1\},
\]
not necessarily with zero diagonal.  If $F=[6]\setminus\{v\}$ and $d_v=J_{vv}$,
then
\[
  \spec(J_{FF})=\{1,1,-d_v,-1,-1\}.
\]
\end{lemma}

\begin{proof}
This follows from Cauchy interlacing.  The first two and last two eigenvalues
are pinned at $1,1$ and $-1,-1$, and the remaining eigenvalue is determined by
the trace.
\end{proof}

\begin{lemma}[Four-by-four blocks on the slice]
Assume $J$ lies on the slice.  If $F=[6]\setminus\{v,w\}$, then
\[
  \spec(J_{FF})=\{1,t,-t,-1\},\qquad t=|a_{vw}|.
\]
\end{lemma}

\begin{proof}
Interlacing pins one eigenvalue at $1$ and one at $-1$.  The trace of $J_{FF}$
is zero on the slice, so the remaining two eigenvalues sum to zero.  Their
squared sum is fixed by the Frobenius norm of $J_{FF}$, which is controlled by
the row-mass identities and equals $2t^2$ for the remaining pair.
\end{proof}

\section{Known obstructions}

The following observations are intended as guardrails for the proof attempt.

\begin{proposition}[Single-pair boundary obstruction; informal statement]\label{prop:boundary}
There is no universal $\delta_0>0$ such that every qualifying pair in
$\St(6,3)$ with
\[
  0<h(\mu_1)+h(\mu_2)-\frac13\le\delta_0
\]
extends to a $1/6$-good triple through that pair.
\end{proposition}

\begin{proof}[Proof status]
This is a recorded exact-arithmetic obstruction in the accompanying notes.  It
is obtained by perturbing a Nesterenko extremal along a Grassmann tangent
direction.  The selected pair moves into positive excess, all four triples
through it move below the threshold to first order, while an external active
triple moves upward and preserves the global GTZ value locally.
\end{proof}

\begin{remark}
Proposition~\ref{prop:boundary} says that Lemma~\ref{lem:extension} is
essentially the maximal reach of any proof using only one pair and its
second-moment data.
\end{remark}

\begin{proposition}[The icosahedral obstruction]\label{prop:ico}
On the equal-leverage icosahedral frame, every pair satisfies
\[
  h(\mu_1)+h(\mu_2)=\frac{13}{11}>\frac13,
\]
so no pair lies in the extension branch of Lemma~\ref{lem:extension}.  However
good triples exist; indeed the coherent triples are precisely the good triples.
\end{proposition}

\begin{proof}
For the icosahedral equiangular tight frame, $\ell_i=1/2$ and
$c_{ij}^2=1/20$ for all $i\ne j$.  Hence each pair has eigenvalues
\[
  \mu_{1,2}=\frac12\pm\frac1{\sqrt{20}},
\]
and a direct substitution gives the displayed $h$-sum.  For a triple, the
off-diagonal signs can be switched to either all positive or all negative
according to the sign of $c_{ij}c_{ik}c_{jk}$.  The all-positive case has
least eigenvalue $1/2-1/\sqrt{20}>1/6$, while the all-negative case has least
eigenvalue $1/2-2/\sqrt{20}<1/6$.
\end{proof}

\section{A minimax/KKT reformulation}

The obstruction above suggests replacing a single-pair view by the active set
of triples for the nonsmooth objective
\[
  F(A)=\max_{|T|=3}\lmin(P_{TT}).
\]
Let $\mathcal A(A)$ denote the active triples at $A$:
\[
  \mathcal A(A)=\{T: \lmin(P_{TT})=F(A)\}.
\]

\begin{definition}
A point $A\in\St(6,3)$ is called a KKT point of $F$ if it is Clarke-stationary
for the maximum of the active eigenvalue functions.  At points where the active
least eigenvalues are simple, this means
\[
  0\in\conv\{\nabla \lmin(P_{TT}):T\in\mathcal A(A)\}
  \subset T_A^*\St(6,3),
\]
with gradients projected to the Stiefel tangent space.
\end{definition}

\begin{proposition}[KKT reduction]\label{prop:kkt}
Conjecture~\ref{conj:gtz63} is equivalent to the assertion that no KKT point of
$F$ has value below $1/6$.
\end{proposition}

\begin{proof}
Since $\St(6,3)$ is compact and $F$ is continuous, $F$ attains a global
minimum.  Every global minimizer is KKT/Clarke-stationary.  Hence if no KKT
point has value below $1/6$, the global minimum is at least $1/6$.  The converse
is immediate.
\end{proof}

\begin{remark}[Local extremals]
Numerical evidence in the notes suggests that the known equality examples are
strict local minima of $F$: their active triple gradients positively span the
tangent dual.  A useful finite task is to turn this into exact certificates,
one equality type at a time.
\end{remark}

\section{A possible route for the proof}

The observations above suggest the following work plan.

\begin{enumerate}[label=\textbf{Step \arabic*.},leftmargin=*]
\item Use Case A and duality to restrict attention to the core region
$\ell_i>1/6$.

\item Prove the equal-leverage slice theorem:
for every zero-diagonal symmetric orthogonal $J$ with spectrum $(+1)^3(-1)^3$,
some triple satisfies the criterion of Lemma~\ref{lem:slicecriterion}.
The current candidate extremal inside the slice is the pentagonal matching
configuration, which appears to have margin about $0.0394$ above $1/6$.

\item Prove exact local KKT certificates at the known equality configurations.
This should be a finite algebraic problem: compute the active gradients and
exhibit an exact positive convex dependence.

\item Enumerate possible active-set types for a hypothetical KKT point with
$F<1/6$ and rule them out, using the pinned block spectra, complementarity, and
the Schur slack identities.

\item If the previous step becomes too large, set up a symmetry-reduced SOS
certificate with the equality manifold imposed as equations.  Low-degree
generic relaxations are expected to be too weak because averaging alone is too
weak on the slice.
\end{enumerate}

\section{Checklist of statements and status}

\begin{center}
\begin{tabular}{p{0.66\linewidth}p{0.22\linewidth}}
\hline
Statement & Current status\\
\hline
Case A reduction & proved\\
Schur criterion and extension lemma & proved\\
Trace branch $\ell_i+\ell_j\ge13/9$ & proved\\
Equal-leverage involution model & proved\\
Slice triple criterion & proved\\
Pinned $5\times5$ and $4\times4$ spectra & proved\\
Single-pair boundary obstruction & proved in accompanying notes\\
Icosahedral obstruction and coherent rescue & proved\\
Equal-leverage slice theorem & open / numerically supported\\
Exact KKT positive spanning at equality examples & numerical evidence\\
No KKT point below $1/6$ & open\\
Full Conjecture~\ref{conj:gtz63} & open\\
\hline
\end{tabular}
\end{center}

\begin{thebibliography}{9}

\bibitem{GTZ1997}
S.~A. Goreinov, E.~E. Tyrtyshnikov, and N.~L. Zamarashkin.
\newblock A theory of pseudoskeleton approximations.
\newblock \emph{Linear Algebra and its Applications}, 261:1--21, 1997.

\bibitem{SenguptaPautov}
R.~Sengupta and M.~Pautov.
\newblock On the submatrices with the best-bounded inverses.
\newblock arXiv:2604.05944, 2026.

\bibitem{Nesterenko}
Y.~Nesterenko.
\newblock Submatrices with the best-bounded inverses: revisiting the hypothesis.
\newblock arXiv preprint.

\end{thebibliography}

\end{document}
